% TOP_PROBLEM: {"id": 5500007, "title": "$k$-sets", "category_id": 2, "category": {"id": 2, "name": "combinatorics", "display_name": "Combinatorics", "description": "Counting problems, graph theory, discrete structures.", "slug": "combinatorics", "order_index": 2, "created_at": "2026-07-31T15:26:25.671Z"}, "status": "open"}
% FIRST_POSTED: 2026-09-25
% ENTRY_KIND: statement_only
% !TeX program = lualatex
% Definition and sources only; this file is not an LLM solution attempt.
\documentclass[11pt]{article}
\usepackage[margin=25mm]{geometry}
\usepackage{fontspec}
\setmainfont{Latin Modern Roman}
\usepackage{amsmath,amssymb,mathtools}
\usepackage{unicode-math}
\setmathfont{Latin Modern Math}
\usepackage{xurl,hyperref}
\hypersetup{colorlinks=true,urlcolor=blue}
\setcounter{secnumdepth}{0}
\setlength{\parindent}{0pt}
\setlength{\parskip}{5pt}
\setlength{\emergencystretch}{3em}
\begin{document}
\section*{\texorpdfstring{$k$-sets}{\$k\$-sets}}\label{5500007}
\subsection{Definitions and mathematical statement}
For an $n$-point set $P\subset{\mathbb{R}}^2$, a $k$-set is a subset $Q\subset P$ with $|Q|=k$ that can be strictly separated from $P\setminus Q$ by a line. Let
\[
 K(n,k)=\max_{|P|=n}\#\{Q\subset P:Q\text{ is a }k\text{-set}
                                      \},\qquad1\le k<n.
\]
Determine its extremal growth as a function of $n$ and $k$, with the general-position convention of the cited problem. In a simple arrangement of nonvertical lines, the $k$-level consists of points on the lines with exactly $k$ lines strictly below, with closures at vertices. Point--line duality relates the two extremal problems, with a possible index shift depending on convention.

The planar case corresponds to line arrangements. Arrangements of hyperplanes in higher dimensions introduce an additional dimension parameter; they are related generalizations, not literally the same planar function.
\par\smallskip\textit{Formulation and concept references:} \href{https://topp.openproblem.net/p7}{[S1]}.

\subsection{Short English statement}
What is the largest possible number of k-point subsets cut off by a line from an n-point planar set, across all values of n and k?
\subsection{Sources}
\begin{itemize}
\item[S1]Formal statement, Problem 7: Statement.
\url{https://topp.openproblem.net/p7}.
\textit{Formulation source. Consulted 24 September 2026: Planar k-set definition and the line-arrangement comparison.}
\end{itemize}
\end{document}
